\section{Related Work}
\label{sec:related}

\paragraph{Agentic Memory}
Much LLM-based research lately has been focused on agents, which use LLMs as the reasoning core of a larger compound system. While agents have shown significant improvements in performance across a variety of domains, a particularly difficult part of working with agents remains managing memory. Due to context length limitations and cost constraints, the agent's entire history cannot be included in every single message. Thus, memories must be encoded, processed, and retrieved only as necessary. Though there are many ways to do this, the most popular way is retrieval-based augmentation (RAG) which projects memories and user queries into a shared embedding space. Memories are then retrieved based on their similarity to the query in this space.

Memories can be as simple as direct copies of inputs and outputs, but frequently they are much more complex. Reflexion \citep{NEURIPS2023_1b44b878} regularly prompts the agent to self-reflect on its actions, and it is these reflections that are stored as memories. In Voyager \citep{Voyager}, memories are agent-created tools that may be re-used as needed. A-MEM \citep{xu2025amem} dynamically organizes memories into interconnected knowledge networks through adaptive tagging and linking, allowing the memory structure to evolve as new information is integrated.

A growing body of work draws on cognitive-architecture abstractions \citep{sumers2024cognitive,huet2025episodic} that separate \emph{episodic} traces of specific past experiences from \emph{semantic} knowledge generalized across them, and explores hierarchical organizations spanning these levels. G-Memory \citep{zhang2025gmemory} structures multi-agent interaction history into three tiers --- insight, query, and interaction graphs --- to support cross-trial knowledge transfer. H-MEM \citep{sun2025hmem} arranges memories across multiple semantic-abstraction layers with index-based routing for efficient retrieval in long-horizon agents, and MemInsight \citep{salama-etal-2025-meminsight} autonomously augments raw interactions with semantic attributes for more contextualized retrieval. MemoryBank \citep{Zhong_Guo_Gao_Ye_Wang_2024} similarly maintains long-term user memories with periodic summarization. Within this landscape, our framework adopts a deliberately simple two-tier structure --- instance-level episodic critiques and a single dataset-level semantic summary --- chosen for reproducibility and interpretability rather than for scaling to long-horizon multi-agent settings. This approach is appropriate for the relatively small datasets in this paper, but future work may benefit from exploring more complex hierarchical memory structures when adapting these techniques to larger datasets.

\paragraph{Fine-tuning}

Though general-purpose large language models can be very capable at certain tasks, they are limited by what is in their training data. To create useful models for niche use-cases, the most obvious approach is to fine-tune a foundation model on a new dataset. Using the principles of transfer learning, it is possible to greatly specialize models to specific tasks \cite{dodge2020finetuningpretrainedlanguagemodels}.

However, this is not without downsides. Though fine-tuning requires much less data than training foundation models from scratch, it still requires a very large amount of labeled data \cite{vieira2024datadatafinetuninglarge}, which may be difficult and expensive to obtain. Changing a model's weights often leads to "catastrophic forgetting," in which the fine-tuned model loses much of its general capabilities \cite{catastrophic_forgetting}. Fine-tuning is more computationally demanding than simple inference, and can be extremely expensive on larger models \cite{hu2022lora}. Lastly, fine-tuning requires access to the model's weights, meaning it cannot be done at all to closed-source models.

\paragraph{In-Context Learning}

Due to these issues with fine-tuning large language models, much research has been done investigating other learning mechanisms that instead treat the models as black boxes. These have come to be referred to generally as in-context learning strategies \citep{dong-etal-2024-survey}.

Largely, they all work on the same concept. LLMs are probabilistic models, and the trajectory of their output depends entirely on the input provided in their context window. By carefully adjusting what input are given to the model, it is possible to shift the output trajectory to something more favorable. This can take the form of presenting new information to the model, or simply changing how the question is asked.

In \citet{ZeroShot}, it was demonstrated that simply appending "Let's think step by step" to a question significantly improved performance across a variety of domains. Safety researchers have determined it is possible to counteract alignment fine-tuning efforts, or "jailbreak" models, through many different kinds of prompts, such as asking questions in code or presenting the question as a matter of life or death.

Beyond simply changing how questions are worded, it is also possible to improve performance by presenting new information. As demonstrated in \citet{FewShot}, presenting several examples of questions and the associated correct answers, or few-shot learning, is an extremely powerful technique to improve performance. 

Reflection-based in-context learning strategies are also becoming popular methods of specializing models to new datasets without relying on fine-tuning \citep{NEURIPS2023_1b44b878, yao2023react, NEURIPS2023_91edff07, yang2025thelight}. Instead of just presenting a few examples of correct question-answer pairs like in few-shot learning, the model makes decisions and then reasons over feedback it receives over the results of those decisions. By iteratively learning the effects of its actions, the model is able to learn how to make better decisions autonomously. However, this research has almost always been focused on giving the model feedback from a simulated environment \citep{Voyager, zhang2025spiralsymbolicllmplanning}, mimicking more traditional reinforcement-learning approaches. Relatively little research has been done on reflection/critique-based in-context learning using other kinds of feedback.
